The errors of the classical surface-area methods are not random, and that is what makes them
correctable. Voxel-face counting is always about 50\% too high, binary marching cubes keeps an error that
does not shrink with resolution, and greyscale marching cubes is pulled down on curved surfaces, by an
amount a simple blurring argument predicts, and pushed up by noise.

A correction learned from synthetic shapes with known area removes most of this error. On test shapes
it cuts the error of greyscale marching cubes from 4.78\% to 0.68\% at $32^3$ and from 3.52\% to 0.18\% at
$64^3$, barely changes with noise, and stays below 2\% on shape families never seen in training. A
Random Forest on ten simple measurements does as well as a 3D CNN with 3.6 million parameters; the CNN
is worth its extra cost mainly for small objects on coarse grids. Structures thinner than about two
voxels and roughness smaller than one voxel remain out of reach, because the information is no longer
in the image. Natural next steps are a more realistic CT simulation, volumes with many touching
structures, and a test on a real scan of an object with a known surface area.
